%======================================================
%  DOCUMENTATION PIPELINE ECOLAB — Overleaf
%  Auteur : Enrique — Mars 2026
%======================================================
\documentclass[12pt,a4paper]{article}

%--- Encodage & langue ---
\usepackage[utf8]{inputenc}
\usepackage[T1]{fontenc}
\usepackage[french]{babel}

%--- Mise en page ---
\usepackage[top=2.5cm, bottom=2.5cm, left=2.8cm, right=2.8cm]{geometry}
\usepackage{setspace}
\setstretch{1.15}

%--- Polices & couleurs ---
\usepackage{lmodern}
\usepackage[dvipsnames,svgnames,table]{xcolor}

%--- Graphiques & figures ---
\usepackage{graphicx}
\usepackage{float}
\usepackage{caption}
\usepackage{subcaption}
\usepackage{tikz}
\usetikzlibrary{shapes.geometric, arrows.meta, positioning, fit, backgrounds, calc}

%--- Tableaux ---
\usepackage{booktabs}
\usepackage{tabularx}
\usepackage{longtable}
\usepackage{array}
\usepackage{multirow}
\usepackage{colortbl}

%--- Listes & puces ---
\usepackage{enumitem}

%--- Liens & PDF ---
\usepackage[colorlinks=true, linkcolor=NavyBlue, citecolor=ForestGreen,
            urlcolor=RoyalBlue, bookmarks=true]{hyperref}
\usepackage{bookmark}

%--- Code source ---
\usepackage{listings}
\lstset{
  basicstyle=\ttfamily\small,
  keywordstyle=\color{NavyBlue}\bfseries,
  commentstyle=\color{gray}\itshape,
  stringstyle=\color{ForestGreen},
  breaklines=true,
  frame=single,
  rulecolor=\color{gray!40},
  backgroundcolor=\color{gray!5},
  numbers=left,
  numberstyle=\tiny\color{gray},
  numbersep=6pt,
  xleftmargin=12pt,
  captionpos=b
}
\lstdefinelanguage{Python3}{
  language=Python,
  morekeywords={True,False,None,async,await,yield,from,import,as,with,assert,lambda}
}

%--- Boîtes colorées ---
\usepackage[most]{tcolorbox}
\tcbuselibrary{skins,breakable}

\newtcolorbox{infobox}[2][]{
  colback=SkyBlue!10, colframe=RoyalBlue!70,
  title=#2, fonttitle=\bfseries\small,
  #1
}
\newtcolorbox{warnbox}[2][]{
  colback=Orange!10, colframe=Orange!80,
  title=#2, fonttitle=\bfseries\small,
  #1
}
\newtcolorbox{codebox}[2][]{
  colback=gray!5, colframe=gray!40,
  title=#2, fonttitle=\ttfamily\small,
  #1
}

%--- En-têtes & pieds de page ---
\usepackage{fancyhdr}
\pagestyle{fancy}
\fancyhf{}
\fancyhead[L]{\small\textit{Pipeline EcoLab v2}}
\fancyhead[R]{\small\textit{\leftmark}}
\fancyfoot[C]{\thepage}
\renewcommand{\headrulewidth}{0.4pt}

%--- Macros utiles ---
\newcommand{\code}[1]{\texttt{\small #1}}
\newcommand{\gpkg}{\textsc{GeoPackage}}
\newcommand{\cosia}{\textsc{COSIA}}
\newcommand{\bdtopo}{\textsc{BD TOPO}}
\newcommand{\scot}{\textsc{SCoT}}
\newcommand{\todo}[1]{\textcolor{red}{\textbf{[TODO: #1]}}}

%--- Couleurs de score ---
\definecolor{scoreQ1}{HTML}{1a7a2e}
\definecolor{scoreQ2}{HTML}{7ec850}
\definecolor{scoreQ3}{HTML}{f5e642}
\definecolor{scoreQ4}{HTML}{f5a623}
\definecolor{scoreQ5}{HTML}{d0021b}

%======================================================
\begin{document}

%--- Page de titre ---
\begin{titlepage}
  \centering
  \vspace*{2cm}
  {\Huge\bfseries Pipeline EcoLab\par}
  \vspace{0.5em}
  {\Large\color{NavyBlue} Documentation technique complète\par}
  \vspace{2cm}
  \rule{\textwidth}{1pt}
  \vspace{1em}
  {\large
  De la modélisation thématique (BERTopic) \\[0.3em]
  à la cartographie interactive \\[0.3em]
  des scores d'artificialisation du territoire SCoT \par
  }
  \vspace{1em}
  \rule{\textwidth}{1pt}
  \vspace{2cm}
  \begin{tabular}{ll}
    \textbf{Auteur :}   & Enrique \\[0.3em]
    \textbf{Version :}  & 2.0 \\[0.3em]
    \textbf{Date :}     & Mars 2026 \\[0.3em]
    \textbf{Territoire :} & SCoT Drôme Ardèche (D007, D026, D084) \\[0.3em]
    \textbf{Corpus :}   & $\approx$50\,000 articles de presse environnementale \\[0.3em]
    \textbf{Données géo :} & $\approx$15 millions de polygones COSIA \\
  \end{tabular}
  \vfill
  {\small Ce document décrit l'intégralité de la chaîne de traitement : scraping, modélisation BERTopic, scoring LLM, analyse géospatiale, export QGIS et application web interactive.}
\end{titlepage}

\tableofcontents
\clearpage

%======================================================
\section{Introduction et vue d'ensemble}

\subsection{Objectif du projet}

La pipeline EcoLab vise à produire une \textbf{carte de score d'artificialisation} à l'échelle parcellaire sur le territoire du \scot{} (Drôme-Ardèche). Pour chaque polygone de la couche \cosia{}, un score compris entre $-2$ (très favorable à la nature) et $+2$ (très favorable à l'artificialisation) est calculé à partir de deux sources complémentaires :

\begin{enumerate}
  \item \textbf{Le signal presse} : des articles de presse environnementale sont analysés par un modèle de langage (Ollama / Mistral) et regroupés par thème via BERTopic. Chaque classe d'occupation du sol hérite du score moyen des thèmes qui lui sont associés.
  \item \textbf{Le contexte géographique} : les couches \bdtopo{} (parcs naturels, forêts publiques, zones de végétation, cours d'eau, barrages, éoliennes…) enrichissent le score via un croisement spatial.
\end{enumerate}

Le score final est une combinaison pondérée 50/50 des deux sources, et est visualisé via QGIS et une application web Streamlit/Leaflet.

\subsection{Vue d'ensemble de la pipeline}

\begin{figure}[H]
\centering
\begin{tikzpicture}[
  node distance=0.8cm and 1.2cm,
  box/.style={rectangle, rounded corners=4pt, draw, minimum width=4.2cm, minimum height=0.85cm, align=center, font=\small},
  data/.style={box, fill=SkyBlue!20, draw=RoyalBlue!60},
  proc/.style={box, fill=ForestGreen!15, draw=ForestGreen!60},
  out/.style={box, fill=Orange!15, draw=Orange!60},
  arr/.style={-{Stealth[length=5pt]}, thick},
  grp/.style={draw=gray!40, dashed, rounded corners=6pt, inner sep=8pt}
]

% Données sources
\node[data] (arts) {Articles presse\\(7 sources CSV)};
\node[data, right=1.5cm of arts] (cosia) {Tuiles COSIA\\(3 dép., $\sim$15M poly.)};
\node[data, right=1.5cm of cosia] (bdtopo) {\bdtopo{} SHP\\(3 dép., 14 catégories)};
\node[data, right=1.5cm of bdtopo] (scot) {Masque SCoT\\(Parcelles\_SCOT.gpkg)};

% BERTopic
\node[proc, below=1.2cm of arts] (bertopic) {\textbf{BERTopic}\\Modélisation thématique};
\draw[arr] (arts) -- (bertopic);

% Étape 01
\node[proc, below=1cm of bertopic] (e01) {\textbf{Étape 01}\\Scoring Ollama/Mistral\\($-2$ à $+2$ par article)};
\draw[arr] (bertopic) -- (e01);

% Étape 03
\node[proc, right=3.5cm of e01] (e03) {\textbf{Étape 03}\\Score COSIA géospatial\\(topic $\rightarrow$ classe)};
\draw[arr] (e01) -- (e03);
\draw[arr] (cosia) -- (e03);
\draw[arr] (scot) -- (e03);

% Étape 04
\node[proc, below=1cm of e03] (e04) {\textbf{Étape 04}\\Enrichissement \bdtopo{}\\(croisements spatiaux)};
\draw[arr] (e03) -- (e04);
\draw[arr] (bdtopo) -- (e04);

% Étape 05
\node[proc, below=1cm of e04] (e05) {\textbf{Étape 05}\\Export QGIS\\(quintiles + QML)};
\draw[arr] (e04) -- (e05);

% Étapes 06
\node[proc, left=2cm of e05] (e06a) {\textbf{Étape 06a}\\Articles par parcelle\\(embeddings + LLM)};
\node[proc, below=0.8cm of e06a] (e06b) {\textbf{Étape 06b}\\Résumés \bdtopo{}\\(combinaisons)};
\node[proc, below=0.8cm of e06b] (e06c) {\textbf{Étape 06c}\\Injection GPKG\\(resume\_bdtopo)};
\draw[arr] (e05) -- (e06a);
\draw[arr] (e06a) -- (e06b);
\draw[arr] (e06b) -- (e06c);

% Output GPKG
\node[out, below=1cm of e06c] (gpkg) {\textbf{GPKG Final}\\COSIA\_BD\_ecolab\_carte\_finale};
\draw[arr] (e06c) -- (gpkg);
\draw[arr] (e05) -- (gpkg);

% App web
\node[out, right=2cm of gpkg] (app) {\textbf{Application web}\\Streamlit + Flask\\+ Leaflet.js};
\draw[arr] (gpkg) -- (app);

\end{tikzpicture}
\caption{Schéma général de la pipeline EcoLab v2}
\end{figure}

\subsection{Volumes de données}

\begin{table}[H]
\centering
\begin{tabular}{lrl}
\toprule
\textbf{Composant} & \textbf{Volume} & \textbf{Remarques} \\
\midrule
Articles de presse & $\sim$50\,000 & Dédupliqués par URL \\
Topics BERTopic & 125+ & Variable selon le run \\
Classes COSIA & 15 & Fixes (nomenclature) \\
Catégories \bdtopo{} & 14 & Fixes (nomenclature) \\
Polygones COSIA (final) & $\sim$15 millions & 3 départements \\
GPKG final & 29,7 Go & Index R-Tree + 15 index SQL \\
Articles synthétiques & 29 & 15 COSIA + 14 \bdtopo{} \\
Combos \bdtopo{} avec résumé & $\sim$160 & Après seuil MIN\_POLYGONS \\
\bottomrule
\end{tabular}
\caption{Volumes de données de la pipeline}
\end{table}

%======================================================
\section{Modélisation thématique — BERTopic}

\subsection{Objectif et principe}

BERTopic est un modèle de \textit{topic modeling} neuronal qui exploite des embeddings de phrases (Sentence-BERT) pour regrouper des articles en thèmes cohérents. Il résout le problème clé suivant : \textbf{comment distinguer automatiquement les articles pertinents pour l'artificialisation des sols de ceux qui ne le sont pas ?}

Le workflow est :
\begin{enumerate}
  \item Encoder chaque article en vecteur dense (modèle multilingue)
  \item Réduire la dimension avec UMAP
  \item Clusteriser avec HDBSCAN
  \item Extraire les mots-clés de chaque cluster via TF-IDF de classe (c-TF-IDF)
  \item Attribuer un label humain-lisible à chaque topic
\end{enumerate}

\subsection{Entraînement (\code{bertopic\_train\_save.py})}

\subsubsection{Sources de données}

\begin{table}[H]
\centering
\begin{tabular}{lll}
\toprule
\textbf{Source} & \textbf{Fichier CSV} & \textbf{Type} \\
\midrule
Cerema & \code{cerema\_actualites.csv} & Actualités institutionnelles \\
Actu-Environnement & \code{actu\_environnement\_actualites.csv} & Presse spécialisée \\
Le Dauphiné & \code{ledauphine\_articles.csv} & Presse régionale \\
Le Monde & \code{lemonde\_articles.csv} & Presse nationale \\
Vert.eco & \code{vert\_eco\_articles.csv} & Presse écologie \\
\bottomrule
\end{tabular}
\caption{Sources d'articles utilisées}
\end{table}

\subsubsection{Paramètres du modèle}

\begin{lstlisting}[language=Python3, caption={Initialisation de BERTopic (bertopic\_train\_save.py)}]
# Modele d'embeddings multilingue
embedding_model = SentenceTransformer(
    "paraphrase-multilingual-MiniLM-L12-v2"
)

# Reduction de dimension
umap_model = UMAP(
    n_components=5,
    metric="cosine",
    random_state=42
)

# Clustering
hdbscan_model = HDBSCAN(
    min_cluster_size=20,
    prediction_data=True
)

# Vectorisation TF-IDF avec stopwords francais
vectorizer = CountVectorizer(
    ngram_range=(1, 2),
    stop_words=FRENCH_STOPWORDS  # ~45 mots
)

# Modele final
topic_model = BERTopic(
    embedding_model=embedding_model,
    umap_model=umap_model,
    hdbscan_model=hdbscan_model,
    vectorizer_model=vectorizer,
    nr_topics=50,
    min_topic_size=20
)
\end{lstlisting}

\subsubsection{Pré-traitement}

\begin{itemize}
  \item Mise en minuscules, normalisation des espaces
  \item Filtrage : documents $<$ 80 caractères supprimés
  \item Déduplication par URL
\end{itemize}

\subsubsection{Sorties}

\begin{itemize}
  \item \code{docs\_with\_topics.csv} — chaque article avec son topic et sa probabilité
  \item \code{topics\_info.csv} — métadonnées (topic ID, nom, top mots-clés, n documents)
  \item \code{bertopic\_model/} — modèle sérialisé (réutilisable)
\end{itemize}

\subsection{Séparation pertinent / non-pertinent}

Après l'entraînement, les topics sont examinés manuellement et mappés :

\begin{itemize}
  \item \textbf{Topics pertinents (125+)} : associés à au moins une classe COSIA ou catégorie \bdtopo{} dans \code{TOPIC\_COSIA\_MAP} et \code{TOPIC\_BDTOPO\_MAP} de \code{config\_v2.py}
  \item \textbf{Topic $-1$} : documents hors-thème (outliers HDBSCAN), ignorés
  \item \textbf{Topics non mappés} : thèmes généraux (politique, santé, sport…) sans lien territorial, ignorés dans le scoring
\end{itemize}

\begin{infobox}{Exemple de mapping topic → classe COSIA}
\small
\code{TOPIC\_COSIA\_MAP = \{}\\
\code{\ \ 3: ["Bâtiment", "Zone imperméable"],\ \ \# topic 3 = urbanisme}\\
\code{\ \ 7: ["Feuillu", "Conifère"],\ \ \ \ \ \ \ \ \# topic 7 = forêt}\\
\code{\ \ 12: ["Surface eau"],\ \ \ \ \ \ \ \ \ \ \ \ \ \ \ \# topic 12 = eau}\\
\code{\ \ ...\}}
\end{infobox}

\subsection{Visualisation (\code{bertopic\_code\_heatmap\_visu.py})}

Plusieurs graphiques sont générés pour explorer et valider le modèle :

\begin{enumerate}
  \item \textbf{Distribution} : histogramme (topic ID vs. nombre de documents)
  \item \textbf{Heatmap topic × source} : crosstab des counts par source de données
  \item \textbf{UMAP documents} : nuage de points 2D coloré par topic (outliers inclus)
  \item \textbf{UMAP topics} : centroïdes des topics avec taille proportionnelle au nombre de documents, labels TF-IDF affichés
  \item \textbf{Tableau récapitulatif} : export CSV (topic ID, nom, n\_docs)
\end{enumerate}

%======================================================
\section{Étape 01 — Scoring Ollama des articles}

\textbf{Script :} \code{01\_topic\_notes\_ollama\_v2.py}

\subsection{Objectif}

Pour chaque article du corpus (après assignation des topics par BERTopic), un score d'artificialisation est attribué par un LLM local (Mistral via Ollama) selon une grille de $-2$ à $+2$.

\subsection{Grille de scoring}

\begin{table}[H]
\centering
\begin{tabular}{>{\centering}p{1.5cm}p{10cm}}
\toprule
\textbf{Score} & \textbf{Signification} \\
\midrule
\rowcolor{scoreQ1!30} $-2$ & Résiste fortement à l'artificialisation (biodiversité, qualité de l'eau, espaces agricoles, zones protégées) \\
\rowcolor{scoreQ2!30} $-1$ & Dominante environnementale légère (mobilité douce, agriculture durable) \\
\rowcolor{scoreQ3!30} $0$  & Pas de lien territorial / neutre \\
\rowcolor{scoreQ4!30} $+1$ & Dominante développement légère (transports, réhabilitation de friches) \\
\rowcolor{scoreQ5!30} $+2$ & Favorise l'artificialisation (construction, infrastructure, immobilier) \\
\bottomrule
\end{tabular}
\caption{Grille de scoring Ollama}
\end{table}

\subsection{Architecture d'appel}

\begin{lstlisting}[language=Python3, caption={Appel Ollama avec parsing robuste du score}]
def ollama_call(prompt: str) -> str:
    """Appel POST a l'API Ollama locale (localhost:11434)."""
    payload = {
        "model": "mistral",
        "prompt": prompt,
        "stream": False
    }
    r = requests.post(OLLAMA_URL, json=payload, timeout=120)
    return r.json()["response"]

def parse_score(raw: str) -> int | None:
    """Extrait un entier entre -2 et +2 de la reponse LLM."""
    matches = re.findall(r"[-+]?\d", raw)
    for m in matches:
        v = int(m)
        if -2 <= v <= 2:
            return v
    return None
\end{lstlisting}

\subsection{Logique de reprise (checkpoint)}

Toutes les 50 articles, le fichier \code{docs\_scored\_v2.csv} est sauvegardé. Au prochain lancement, les articles déjà scorés sont sautés automatiquement — le traitement peut être interrompu et repris.

\subsection{Sorties}

\begin{itemize}
  \item \code{docs\_scored\_v2.csv} — tous les articles avec leur score LLM
  \item \code{topic\_score\_summary\_v2.csv} — score moyen/médian par topic (agrégation par \code{groupby("topic")})
\end{itemize}

%======================================================
\section{Étape 02 — Visualisation des topics}

\textbf{Script :} \code{02\_topic\_viz\_v2.py}

\subsection{Graphiques produits}

\begin{enumerate}
  \item \textbf{Distribution} : barres horizontales, nombre de documents par topic (trié croissant)
  \item \textbf{Scores} : barres horizontales, score moyen Ollama par topic avec codage couleur par quintile :
\end{enumerate}

\begin{table}[H]
\centering
\begin{tabular}{lll}
\toprule
\textbf{Couleur} & \textbf{Hex} & \textbf{Seuil} \\
\midrule
\cellcolor{scoreQ1!60}Rouge foncé & \code{\#d73027} & $\leq -0.5$ (très protecteur) \\
\cellcolor{scoreQ4!40}Orange & \code{\#fc8d59} & $-0.5$ à $0$ (protecteur) \\
\cellcolor{scoreQ3!60}Gris & \code{\#d9d9d9} & $0$ (neutre) \\
\cellcolor{scoreQ2!40}Vert clair & \code{\#91cf60} & $0$ à $0.5$ (développement) \\
\cellcolor{scoreQ1!40}Vert foncé & \code{\#1a9641} & $\geq 0.5$ (très développement) \\
\bottomrule
\end{tabular}
\caption{Codage couleur des scores de topic}
\end{table}

%======================================================
\section{Étape 03 — Score géospatial COSIA}

\textbf{Script :} \code{03\_geospatial\_score\_cosia\_v2.py}

\subsection{Objectif}

Assigner un score d'artificialisation à chaque polygone de la couche \cosia{} en fonction de sa classe d'occupation du sol, via le mapping \code{TOPIC\_COSIA\_MAP}.

\subsection{Construction du score par classe}

Pour chaque classe COSIA :
$$
\text{score\_classe}(c) = \frac{1}{|T_c|} \sum_{t \in T_c} \text{score\_moyen}(t)
$$
où $T_c$ est l'ensemble des topics mappés à la classe $c$.

\subsection{Architecture de traitement parallèle}

\begin{lstlisting}[language=Python3, caption={Lecture parallèle des tuiles COSIA}]
# 6 workers I/O en parallele, traitement sequentiel
with ThreadPoolExecutor(max_workers=6) as executor:
    futures = {
        executor.submit(read_and_filter_tile, tile): tile
        for tile in all_tiles
    }
    for future in as_completed(futures):
        gdf = future.result()
        if gdf is not None and len(gdf) > 0:
            process_gdf(gdf)  # sequentiel pour la memoire
\end{lstlisting}

\subsection{Filtre géographique}

Plutôt qu'un découpage géométrique coûteux, un simple filtre par \textit{bounding box} (Lambert 93) est appliqué : les polygones dont le centroïde est en dehors du SCoT sont exclus. La bbox du SCoT est :
$$
(816\,777,\ 6\,340\,562,\ 914\,455,\ 6\,407\,100) \text{ [Lambert 93]}
$$

\subsection{Sorties}

\begin{itemize}
  \item \code{COSIA\_SCORE\_v2.gpkg} — polygones scorés (\textasciitilde{}15M lignes)
  \item \code{COSIA\_SCORE\_v2\_attrs.csv} — attributs seuls (sans géométrie)
\end{itemize}

%======================================================
\section{Étape 04 — Enrichissement BD TOPO}

\textbf{Script :} \code{04\_enrich\_cosia\_bdtopo\_v2.py}

\subsection{Objectif}

Enrichir chaque polygone COSIA avec son contexte géographique \bdtopo{} : 14 colonnes booléennes \code{bdtopo\_*} indiquent si le polygone intersecte une zone protégée, une forêt publique, un cours d'eau, etc.

\subsection{Les 14 catégories BD TOPO}

\begin{table}[H]
\centering
\begin{tabular}{lll}
\toprule
\textbf{Catégorie} & \textbf{Colonne} & \textbf{Type de zone} \\
\midrule
Natura 2000 & \code{bdtopo\_NATURA2000} & Protection européenne \\
Parc naturel régional & \code{bdtopo\_PARC\_PNR} & Protection nationale \\
Réserve naturelle & \code{bdtopo\_RESERVE\_NAT} & Protection nationale \\
Géoparc & \code{bdtopo\_GEOPARC} & Label UNESCO \\
Forêt publique & \code{bdtopo\_FORET\_PUBLIQUE} & ONF \\
Zone végétation forêt & \code{bdtopo\_ZONE\_VEG\_FORET} & Zone de végétation \\
Zone végétation vigne & \code{bdtopo\_ZONE\_VEG\_VIGNE} & Zone de végétation \\
Zone végétation verger & \code{bdtopo\_ZONE\_VEG\_VERGER} & Zone de végétation \\
Zone végétation lande & \code{bdtopo\_ZONE\_VEG\_LANDE} & Zone de végétation \\
Cours d'eau & \code{bdtopo\_COURS\_EAU} & Hydrographie \\
Surface en eau & \code{bdtopo\_SURFACE\_EAU} & Hydrographie \\
Barrage & \code{bdtopo\_BARRAGE} & Infrastructure \\
Éolienne & \code{bdtopo\_EOLIENNE} & Infrastructure \\
Cours d'eau nommé & \code{nom\_cours\_eau} & (texte, TOPONYME) \\
\bottomrule
\end{tabular}
\caption{Catégories BD TOPO et colonnes correspondantes}
\end{table}

\subsection{Score BD TOPO et score final}

Le score BD TOPO d'un polygone est la moyenne des scores des catégories actives :
$$
\text{score\_bdtopo}(p) = \frac{1}{|C_p|} \sum_{c \in C_p} \text{score\_classe\_bdtopo}(c)
$$

Le score final est une pondération 50/50 :
$$
\text{score\_final}(p) = \begin{cases}
  0.5 \times \text{score\_cosia}(p) + 0.5 \times \text{score\_bdtopo}(p) & \text{si } |C_p| > 0 \\
  \text{score\_cosia}(p) & \text{sinon}
\end{cases}
$$

\subsection{Optimisations}

\begin{itemize}
  \item \textbf{Index spatial R-Tree} construit une seule fois pour tout le GPKG COSIA
  \item \textbf{ZONE\_DE\_VEGETATION} (594\,000 entités) mise en cache après filtrage bbox initial
  \item Opérations vectorisées NumPy pour le calcul du score final
\end{itemize}

\subsection{Sorties}

\begin{itemize}
  \item \code{COSIA\_SCORE\_finale\_ecolab.gpkg} — polygones finaux avec toutes les colonnes
  \item \code{COSIA\_SCORE\_finale\_attrs.csv} — attributs seuls
\end{itemize}

%======================================================
\section{Étape 05 — Export QGIS}

\textbf{Script :} \code{05\_export\_qgis\_v2.py}

\subsection{Quintiles et couleurs}

Le score final est discrétisé en 5 classes par quintiles (20\%, 40\%, 60\%, 80\%) :

\begin{table}[H]
\centering
\begin{tabular}{llll}
\toprule
\textbf{Quintile} & \textbf{Couleur} & \textbf{Hex} & \textbf{Signification} \\
\midrule
\rowcolor{scoreQ1!20} Q1 (0–20\%) & Vert foncé & \code{\#1a7a2e} & Très protecteur \\
\rowcolor{scoreQ2!20} Q2 (20–40\%) & Vert clair & \code{\#7ec850} & Protecteur \\
\rowcolor{scoreQ3!20} Q3 (40–60\%) & Jaune & \code{\#f5e642} & Neutre/mixte \\
\rowcolor{scoreQ4!20} Q4 (60–80\%) & Orange & \code{\#f5a623} & Développement \\
\rowcolor{scoreQ5!20} Q5 (80–100\%) & Rouge & \code{\#d0021b} & Très développement \\
\bottomrule
\end{tabular}
\caption{Quintiles et palette de couleurs}
\end{table}

\subsection{Fichier QML}

Un fichier de style \code{.qml} est généré automatiquement avec les intervalles de quintiles, de sorte que QGIS charge le style dès l'ouverture du \gpkg{}.

\subsection{Grille de tuiles 1km × 1km}

En plus des polygones individuels, une grille de tuiles d'1 km² est calculée avec le score moyen pondéré par surface :
$$
\text{score\_tuile} = \frac{\sum_i \text{score}_i \times \text{aire}_i}{\sum_i \text{aire}_i}
$$

\subsection{Sorties}

\begin{itemize}
  \item \code{COSIA\_BD\_ecolab\_carte\_finale.gpkg} — \gpkg{} final (renommage sans copie)
  \item \code{COSIA\_BD\_ecolab\_carte\_finale.qml} — style QGIS auto-chargé
  \item Couche \code{TUILES\_SCORE} ajoutée au même \gpkg{}
\end{itemize}

%======================================================
\section{Étape 06a — Articles par parcelle}

\textbf{Script :} \code{06\_articles\_par\_parcelle\_v2.py}

\subsection{Objectif}

Associer 3 articles de presse pertinents à chaque classe COSIA et catégorie \bdtopo{}, via une recherche sémantique en deux temps.

\subsection{Workflow en 2 temps}

\begin{enumerate}
  \item \textbf{Recherche par embeddings (cosine similarity)} : Pour chaque label (classe/catégorie), un article synthétique de référence est encodé, puis les 15 articles les plus proches du corpus sont récupérés.
  \item \textbf{Scoring LLM} : Ollama/Mistral note chaque candidat de 0.0 à 1.0 pour sa pertinence. Les 3 meilleurs sont conservés.
\end{enumerate}

\subsection{Articles synthétiques}

Pour chaque label, un article de référence de 300–500 mots est rédigé manuellement, mentionnant explicitement le contexte du SCoT Drôme-Ardèche. Ceci guide la recherche sémantique vers des articles localement pertinents.

\subsection{Cache d'embeddings}

\begin{lstlisting}[language=Python3, caption={Système de cache d'embeddings}]
# Sauvegarde
np.save("articles_embeddings_cache.npy", embeddings_array)
pd.DataFrame({"url": urls}).to_csv("articles_embeddings_index.csv")

# Chargement (si cache existant)
if cache_exists:
    embeddings = np.load("articles_embeddings_cache.npy")
    urls = pd.read_csv("articles_embeddings_index.csv")["url"].tolist()
\end{lstlisting}

\subsection{Résumé COSIA par Ollama}

Pour chaque classe COSIA, un résumé de 3–5 phrases est généré par Mistral à partir des 3 articles retenus. Ce résumé est stocké dans \code{COSIA\_articles\_par\_classe.csv}.

\subsection{Sorties}

\begin{itemize}
  \item \code{COSIA\_articles\_par\_classe.csv} — 15 lignes (une par classe), avec 3 URLs + titres COSIA, 6 URLs + titres contexte, résumé
  \item \code{COSIA\_BD\_ecolab\_articles.qml} — Map Tips QGIS (HTML avec liens cliquables)
\end{itemize}

%======================================================
\section{Étape 06b — Résumés BD TOPO}

\textbf{Script :} \code{06b\_resume\_bdtopo.py}

\subsection{Objectif}

Générer un résumé textuel pour chaque combinaison unique de catégories \bdtopo{} présente dans le \gpkg{}. Par exemple, un polygone situé à la fois en forêt publique et en Natura 2000 reçoit un résumé spécifique à ce contexte.

\subsection{Clé de combinaison}

La clé est une liste triée et séparée par \code{|} des catégories actives :
\begin{center}
\code{"FORET\_PUBLIQUE|NATURA2000|ZONE\_VEG\_FORET"}
\end{center}

\subsection{Seuil et filtrage}

Seules les combinaisons comptant au moins \code{MIN\_POLYGONS} polygones reçoivent un résumé (évite les résumés pour des combinaisons anecdotiques). Ce seuil a été baissé lors du second run pour couvrir des combinaisons rares comme \code{EOLIENNE|FORET\_PUBLIQUE|ZONE\_VEG\_FORET}.

\subsection{Sorties}

\begin{itemize}
  \item \code{BDTOPO\_resume\_par\_combo.csv} — 160 lignes : \code{combo\_key}, \code{n\_polygones}, \code{resume\_bdtopo}
\end{itemize}

%======================================================
\section{Étape 06c — Injection des résumés dans le GPKG}

\textbf{Script :} \code{06c\_inject\_resume\_bdtopo.py}

\subsection{Objectif}

Injecter la colonne \code{resume\_bdtopo} directement dans la table SQLite du \gpkg{} pour éviter une jointure à l'exécution (plus rapide pour QGIS et l'application web).

\subsection{Optimisations SQLite}

\begin{lstlisting}[language=Python3, caption={Pragmas SQLite pour writes rapides}]
conn.execute("PRAGMA journal_mode=MEMORY")  # journal en RAM
conn.execute("PRAGMA synchronous=OFF")      # pas de fsync
conn.execute("PRAGMA cache_size=100000")    # 100K pages en cache
\end{lstlisting}

\subsection{Logique de mise à jour}

Pour chaque combinaison unique :
\begin{lstlisting}[language=Python3, caption={UPDATE ciblé par combo BD TOPO}]
# Construction de la clause WHERE
where_parts = []
for cat in ALL_BDTOPO_CATS:
    val = 1 if cat in active_cats else 0
    where_parts.append(f"bdtopo_{cat} = {val}")

sql = f"""
    UPDATE COSIA_SCORE_finale
    SET resume_bdtopo = ?
    WHERE {' AND '.join(where_parts)}
"""
conn.execute(sql, (resume_text,))
\end{lstlisting}

%======================================================
\section{Application web interactive}

\textbf{Script :} \code{app\_carte2.py}

\subsection{Architecture technique}

L'application combine trois couches :

\begin{figure}[H]
\centering
\begin{tikzpicture}[
  node distance=1cm,
  layer/.style={rectangle, rounded corners=6pt, draw, minimum width=9cm, minimum height=1.2cm, align=center, font=\small},
  arr/.style={-{Stealth[length=5pt]}, thick}
]
\node[layer, fill=RoyalBlue!20, draw=RoyalBlue] (front) {\textbf{Frontend — Leaflet.js (JavaScript)}\\ Rendu cartographique, pan/zoom, sélection de parcelle};
\node[layer, fill=ForestGreen!15, draw=ForestGreen, below=0.7cm of front] (mid) {\textbf{Mid-tier — Flask (Python, port 5050)}\\ API REST : GeoJSON par bbox, détail par FID, filtres};
\node[layer, fill=Orange!15, draw=Orange, below=0.7cm of mid] (back) {\textbf{Backend — SQLite3 + Index R-Tree}\\ GPKG 29.7 Go, 15M polygones, 15 index SQL};
\node[layer, fill=Purple!10, draw=Purple!60, left=0.3cm of front, minimum width=3cm, xshift=-4.5cm] (st) {\textbf{Streamlit}\\ Sidebar filtres\\ Panel détail};

\draw[arr] (front) -- (mid) node[midway, right] {\small HTTP};
\draw[arr] (mid) -- (back) node[midway, right] {\small SQLite};
\draw[arr, <->] (st) -- (front) node[midway, above] {\small iframe};
\end{tikzpicture}
\caption{Architecture de l'application web}
\end{figure}

\subsection{Flask — Routes API}

\begin{table}[H]
\centering
\begin{tabular}{lll}
\toprule
\textbf{Route} & \textbf{Méthode} & \textbf{Rôle} \\
\midrule
\code{/api/features} & GET & GeoJSON des polygones dans la bbox courante \\
\code{/api/detail/<fid>} & GET & Tous les attributs d'un polygone (articles, scores…) \\
\code{/api/filter\_bbox} & GET & Bbox de l'ensemble des résultats filtrés \\
\code{/api/select/<fid>} & GET & Stocker le FID sélectionné en session \\
\bottomrule
\end{tabular}
\caption{Routes de l'API Flask}
\end{table}

\subsection{Filtre et chargement}

\begin{lstlisting}[language=Python3, caption={Requête SQL avec filtres dynamiques (api\_features)}]
conditions = ["1=1"]
params = []

# Filtre par classes COSIA
if classes_filter:
    placeholders = ",".join("?" * len(classes_filter))
    conditions.append(f"classe IN ({placeholders})")
    params.extend(classes_filter)

# Filtres BD TOPO (AND logique)
for cat in bdtopo_filter:
    conditions.append(f"bdtopo_{cat} = 1")

# Filtre bbox viewport
conditions.append("""
    (minx <= ? AND maxx >= ? AND miny <= ? AND maxy >= ?)
""")
params.extend([xmax_l93, xmin_l93, ymax_l93, ymin_l93])

sql = f"""
    SELECT fid, geom, classe, score_cosia
    FROM COSIA_SCORE_finale
    WHERE {' AND '.join(conditions)}
    LIMIT 5000
"""
\end{lstlisting}

\subsection{Système de cache des bbox}

Au démarrage de l'application, un thread de fond calcule et met en cache les \textit{bounding boxes} de chaque combinaison filtre :

\begin{lstlisting}[language=Python3, caption={Pré-calcul de la bbox par classe (background thread)}]
sql = """
    SELECT classe,
           MIN(minx) as xmin, MAX(maxx) as xmax,
           MIN(miny) as ymin, MAX(maxy) as ymax
    FROM COSIA_SCORE_finale
    GROUP BY classe
"""
# Resultat sauvegarde dans bbox_cache.json
# → zoom automatique instantane au filtrage suivant
\end{lstlisting}

\subsection{Sidebar Streamlit — Filtres}

\begin{figure}[H]
\centering
\begin{tikzpicture}[
  box/.style={rectangle, rounded corners=3pt, draw=gray!60, fill=white, minimum width=5cm, minimum height=0.5cm, font=\small, align=center},
  section/.style={rectangle, rounded corners=3pt, fill=NavyBlue!10, draw=NavyBlue!40, minimum width=5cm, minimum height=0.4cm, font=\small\bfseries, align=center}
]
\node[section] (s1) {Classes COSIA};
\node[box, below=0.1cm of s1] (c1) {Bâtiment};
\node[box, below=0.05cm of c1] (c2) {Zone imperméable};
\node[box, below=0.05cm of c2] (c3) {Surface eau};
\node[box, below=0.05cm of c3] (c4) {... (15 classes)};
\node[section, below=0.3cm of c4] (s2) {Contexte BD TOPO};
\node[box, below=0.1cm of s2] (b1) {Natura 2000};
\node[box, below=0.05cm of b1] (b2) {Parc naturel régional};
\node[box, below=0.05cm of b2] (b3) {... (14 catégories)};
\node[box, below=0.4cm of b3, fill=ForestGreen!15, draw=ForestGreen] (btn) {\textbf{Zoom sur sélection}};
\end{tikzpicture}
\caption{Sidebar Streamlit avec filtres}
\end{figure}

\subsection{Panel de détail parcelle}

Quand l'utilisateur clique sur une parcelle sur la carte :

\begin{enumerate}
  \item Leaflet envoie le FID à Flask (\code{/api/select/<fid>})
  \item Streamlit interroge \code{/api/detail/<fid>} (polling toutes les 2s)
  \item Le panel affiche :
    \begin{itemize}
      \item \textbf{Score global} avec jauge colorée
      \item Score COSIA / Score \bdtopo{} séparément
      \item \textbf{Badges} des catégories \bdtopo{} actives
      \item \textbf{Articles COSIA} (3 liens cliquables) + résumé Mistral
      \item \textbf{Articles contexte} (3–6 liens) pour le contexte \bdtopo{}
      \item \textbf{Résumé \bdtopo{}} généré par Ollama pour la combinaison
    \end{itemize}
\end{enumerate}

\subsection{Carte Leaflet}

\begin{lstlisting}[language=Python3, caption={Seuils de score et couleurs (app\_carte2.py)}]
# Seuils quantiles du GPKG final
SCORE_THRESHOLDS = [-0.531, -0.465, -0.396, -0.297]

# Palette de couleurs
SCORE_COLORS = ["#1a7a2e", "#7ec850", "#f5e642", "#f5a623", "#d0021b"]

# Rectangle SCoT (WGS84) affiche en tirets bleus
SCOT_BOUNDS = [[44.154, 4.460], [44.730, 5.709]]
L.rectangle(SCOT_BOUNDS, {color:'#1976D2', dashArray:'8 5', weight:2})
\end{lstlisting}

\textbf{Threshold de zoom} : les polygones individuels ne se chargent qu'à partir du zoom 13 (environ commune), pour éviter de saturer le navigateur.

\subsection{Performances}

\begin{table}[H]
\centering
\begin{tabular}{ll}
\toprule
\textbf{Opération} & \textbf{Temps typique} \\
\midrule
Requête bbox + GeoJSON (5000 poly.) & $<$ 200 ms \\
Requête détail (1 parcelle) & $<$ 50 ms \\
Calcul bbox filtre (cache chaud) & $<$ 10 ms \\
Calcul bbox filtre (sans cache) & $\sim$ 2–5 s \\
Démarrage Flask + index warmup & $\sim$ 5 s \\
\bottomrule
\end{tabular}
\caption{Performances de l'application web}
\end{table}

%======================================================
\section{Index SQLite et optimisations}

\textbf{Script :} \code{create\_indexes.py}

\subsection{Index créés}

15 index SQLite sont créés sur la table \code{COSIA\_SCORE\_finale} pour accélérer les filtres :

\begin{lstlisting}[language=Python3, caption={Création des index SQLite}]
indexes = [
    ("idx_classe",          "classe"),
    ("idx_bdtopo_NATURA2000", "bdtopo_NATURA2000"),
    ("idx_bdtopo_PARC_PNR",   "bdtopo_PARC_PNR"),
    # ... 12 autres categories bdtopo
]

for idx_name, col_name in indexes:
    conn.execute(f"""
        CREATE INDEX IF NOT EXISTS {idx_name}
        ON COSIA_SCORE_finale ({col_name})
    """)
\end{lstlisting}

\begin{warnbox}{Temps de création}
Chaque index prend $\sim$40 minutes sur le \gpkg{} de 29.7 Go (15M lignes). Au total : $\sim$10 heures. À n'effectuer qu'une seule fois.
\end{warnbox}

%======================================================
\section{Configuration centrale}

\textbf{Fichier :} \code{config\_v2.py}

Toute la pipeline est configurée depuis un seul fichier. Les paramètres principaux :

\begin{lstlisting}[language=Python3, caption={Extraits de config\_v2.py}]
# Chemins
BASE_DIR = Path(__file__).parent
OUT_DIR  = BASE_DIR / "outputs"
GEO_DIR  = OUT_DIR / "geo"

# Ollama
OLLAMA_URL     = "http://localhost:11434/api/generate"
OLLAMA_MODEL   = "mistral"
OLLAMA_TIMEOUT = 120
SLEEP_BETWEEN  = 0.3  # secondes entre les appels

# Departements
COSIA_DEPTS = ["D007", "D026", "D084"]
BDTOPO_DEPTS = ["D007", "D026", "D084"]

# Bbox SCoT [Lambert 93]
SCOT_BBOX = (816777, 6340562, 914455, 6407100)

# Mapping topic -> classes COSIA (125+ topics)
TOPIC_COSIA_MAP = {
    3:  ["Batiment", "Zone impermeable"],
    7:  ["Feuillu", "Conifere"],
    12: ["Surface eau"],
    ...
}

# Mapping topic -> categories BD TOPO (40+ topics)
TOPIC_BDTOPO_MAP = {
    5:  ["NATURA2000"],
    9:  ["FORET_PUBLIQUE", "ZONE_VEG_FORET"],
    ...
}
\end{lstlisting}

%======================================================
\section{Lancement de la pipeline}

\subsection{Prérequis}

\begin{itemize}
  \item Python 3.10+, GDAL/GeoPandas, PyProj, Sentence-Transformers, BERTopic
  \item Ollama installé localement avec le modèle Mistral (\code{ollama pull mistral})
  \item QGIS 3.40+ pour la visualisation
  \item Streamlit (\code{pip install streamlit flask}) pour l'application web
\end{itemize}

\subsection{Ordre d'exécution}

\begin{table}[H]
\centering
\begin{tabular}{lll}
\toprule
\textbf{Ordre} & \textbf{Script} & \textbf{Durée approx.} \\
\midrule
0 & \code{bertopic\_train\_save.py} & 30–60 min \\
1 & \code{01\_topic\_notes\_ollama\_v2.py} & 3–8h (50K articles) \\
2 & \code{02\_topic\_viz\_v2.py} & 1–2 min \\
3 & \code{03\_geospatial\_score\_cosia\_v2.py} & 30 min \\
4 & \code{04\_enrich\_cosia\_bdtopo\_v2.py} & 45–90 min \\
5 & \code{05\_export\_qgis\_v2.py} & 10–20 min \\
6a & \code{06\_articles\_par\_parcelle\_v2.py} & 30–60 min \\
6b & \code{06b\_resume\_bdtopo.py} & 30–60 min \\
6c & \code{06c\_inject\_resume\_bdtopo.py} & 10–30 min \\
— & \code{create\_indexes.py} & $\sim$10h (une fois) \\
\bottomrule
\end{tabular}
\caption{Ordre et durée d'exécution des scripts}
\end{table}

\subsection{Lancement de l'application web}

\begin{lstlisting}[language=bash, caption={Démarrage de l'application Streamlit}]
cd pipeline_v2
streamlit run app_carte2.py
# Interface disponible sur http://localhost:8501
# API Flask sur http://localhost:5050
\end{lstlisting}

%======================================================
\section{Flux de données complet}

\begin{figure}[H]
\centering
\begin{tikzpicture}[
  node distance=0.5cm and 0.3cm,
  file/.style={rectangle, rounded corners=2pt, draw=gray, fill=gray!10, font=\tiny, minimum height=0.45cm, text width=3.2cm, align=center},
  step/.style={rectangle, rounded corners=4pt, draw=NavyBlue!60, fill=NavyBlue!10, font=\small\bfseries, minimum height=0.7cm, text width=3.8cm, align=center},
  arr/.style={-{Stealth[length=4pt]}, gray},
  thickarr/.style={-{Stealth[length=5pt]}, NavyBlue, thick}
]

\node[file] (f0) {docs\_with\_topics.csv};
\node[step, below=0.3cm of f0] (s01) {01 — Scoring Ollama};
\node[file, below=0.3cm of s01] (f1) {docs\_scored\_v2.csv\\topic\_score\_summary\_v2.csv};
\node[step, below=0.3cm of f1] (s03) {03 — Score COSIA géospatial};
\node[file, below=0.3cm of s03] (f2) {COSIA\_SCORE\_v2.gpkg};
\node[step, below=0.3cm of f2] (s04) {04 — Enrichissement BD TOPO};
\node[file, below=0.3cm of s04] (f3) {COSIA\_SCORE\_finale\_ecolab.gpkg};
\node[step, below=0.3cm of f3] (s05) {05 — Export QGIS};
\node[file, below=0.3cm of s05] (f4) {COSIA\_BD\_ecolab\_carte\_finale.gpkg\\ + .qml};
\node[step, below=0.3cm of f4] (s06a) {06a — Articles par classe};
\node[file, below=0.3cm of s06a] (f5) {COSIA\_articles\_par\_classe.csv};
\node[step, below=0.3cm of f5] (s06b) {06b — Résumés BD TOPO combos};
\node[file, below=0.3cm of s06b] (f6) {BDTOPO\_resume\_par\_combo.csv};
\node[step, below=0.3cm of f6] (s06c) {06c — Injection GPKG};
\node[file, below=0.3cm of s06c] (f7) {\textbf{GPKG FINAL}\\(+ resume\_bdtopo)};

\draw[thickarr] (f0) -- (s01);
\draw[thickarr] (s01) -- (f1);
\draw[thickarr] (f1) -- (s03);
\draw[thickarr] (s03) -- (f2);
\draw[thickarr] (f2) -- (s04);
\draw[thickarr] (s04) -- (f3);
\draw[thickarr] (f3) -- (s05);
\draw[thickarr] (s05) -- (f4);
\draw[thickarr] (f4) -- (s06a);
\draw[thickarr] (s06a) -- (f5);
\draw[thickarr] (f5) -- (s06b);
\draw[thickarr] (s06b) -- (f6);
\draw[thickarr] (f6) -- (s06c);
\draw[thickarr] (s06c) -- (f7);

\end{tikzpicture}
\caption{Chaîne de fichiers de bout en bout}
\end{figure}

%======================================================
\section{Résultats et visualisation QGIS}

\subsection{Couches disponibles dans le GPKG final}

\begin{itemize}
  \item \textbf{COSIA\_SCORE\_finale} ($\sim$15M polygones) avec :
    \begin{itemize}
      \item \code{classe} — classe COSIA (15 valeurs)
      \item \code{score\_cosia} — score issu des topics
      \item \code{score\_bdtopo} — score issu du contexte géographique
      \item \code{score\_polygone} — score final (blended 50/50)
      \item 14 colonnes \code{bdtopo\_*} (booléens 0/1)
      \item \code{nom\_cours\_eau} — toponymie (si applicable)
      \item \code{resume\_bdtopo} — résumé Ollama de la combinaison \bdtopo{}
    \end{itemize}
  \item \textbf{TUILES\_SCORE} — grille 1km² avec score moyen pondéré
\end{itemize}

\subsection{Jointure des résumés dans QGIS}

Pour joindre \code{COSIA\_articles\_par\_classe.csv} à la couche vecteur :
\begin{enumerate}
  \item Propriétés de la couche → Jointures → Ajouter
  \item Couche jointe : \code{COSIA\_articles\_par\_classe.csv}
  \item Champ de jointure de la couche : \code{classe}
  \item Champ cible : \code{classe}
\end{enumerate}

%======================================================
\section{Choix techniques et justifications}

\begin{table}[H]
\centering
\begin{tabular}{p{4cm}p{4cm}p{5cm}}
\toprule
\textbf{Choix} & \textbf{Alternative} & \textbf{Justification} \\
\midrule
\gpkg{} (SQLite) & PostGIS & Pas de serveur requis, portable, R-Tree intégré \\
Ollama local & API OpenAI & Pas de coût, données locales, reproductible \\
Leaflet.js & MapLibre GL & Stabilité, facilité d'intégration dans iframe \\
BERTopic & LDA/NMF & Meilleure cohérence des topics, support multilingue \\
Flask thread & FastAPI & Minimal, pas de processus externe supplémentaire \\
CSV pour articles & 15M UPDATE & Léger, performant, jointure côté app \\
Bbox filter & Clip géométrique & 100× plus rapide sur 15M polygones \\
ThreadPool I/O & ProcessPool & Évite la sérialisation GeoDataFrame, GIL ok pour I/O \\
\bottomrule
\end{tabular}
\caption{Décisions d'architecture}
\end{table}

%======================================================
\section{Glossaire}

\begin{description}[leftmargin=3cm, style=nextline]
  \item[\cosia{}] Classification de l'Occupation des Sols par Imagerie Aérienne — nomenclature nationale à 15 classes pour l'occupation du sol.
  \item[\bdtopo{}] Base de Données Topographique de l'IGN — données vectorielles de référence (forêts, cours d'eau, bâtiments…).
  \item[\scot{}] Schéma de Cohérence Territoriale — document de planification à l'échelle intercommunale.
  \item[BERTopic] Modèle de topic modeling neuronal basé sur BERT + HDBSCAN + c-TF-IDF.
  \item[Ollama] Framework local pour faire tourner des LLMs (ici Mistral) sans API payante.
  \item[HDBSCAN] Hierarchical Density-Based Spatial Clustering — algorithme de clustering sans nombre de clusters fixe.
  \item[UMAP] Uniform Manifold Approximation and Projection — réduction de dimension non-linéaire.
  \item[Lambert 93] Système de projection officiel français (EPSG:2154), utilisé dans le \gpkg{}.
  \item[WGS84] Système géodésique mondial (EPSG:4326), utilisé dans l'application web (Leaflet).
  \item[R-Tree] Index spatial en arbre, utilisé par SQLite pour les requêtes géographiques rapides.
  \item[\gpkg{}] GeoPackage — format de base de données géospatiale basé sur SQLite.
  \item[c-TF-IDF] Class TF-IDF — variante de TF-IDF calculée au niveau de chaque classe (topic) plutôt que de chaque document.
\end{description}

%======================================================
\end{document}
